\section{Introduction}\label{sec:introduction}

A group is \emph{periodic} if each of its elements has finite order.
It is \emph{residually finite} if, for every nonidentity element $g$,
there is a homomorphism to a finite group whose value at $g$ is
nonidentity. We use \emph{2-group} to mean a group in which every
element has finite 2-power order; the group itself need not be finite.
An \emph{ordinary finite presentation} is a quotient of a free group
on finitely many generators by the normal closure of finitely many
words. We prove the following.

\begin{theorem}\label{thm:main}
There exists an infinite residually finite group with an ordinary finite
presentation in which every element has finite 2-power order.
\end{theorem}

The group is an abstract discrete group. Finite presentation is meant
in that category, rather than in a variety of groups or in the category
of pro-2 groups. The orders of its elements are necessarily unbounded.

\subsection*{History and relation to earlier constructions}

Burnside's finiteness questions \cite{Burnside1902} ask how finite-order
conditions constrain a group with finitely many generators. Golod's
construction \cite[Example~2]{Golod1964}, based on the
Golod--Shafarevich dimension argument \cite{GolodShafarevich1964},
already produces infinite finitely generated residually finite
$p$-groups. Its algebraic mechanism is particularly relevant here:
nilpotence gives finite $p$-power orders, while a grading separates
elements in finite quotients. Grigorchuk later gave a concrete infinite
finitely generated residually finite 2-group
\cite{Grigorchuk1980}; his original theorem also establishes that this
group is not finitely presented. Thus residual finiteness and fixed-prime
torsion can coexist with finite generation. The additional requirement
in Theorem~\ref{thm:main} is an ordinary finite presentation.

The question whether an infinite periodic group can have an ordinary
finite presentation is recorded in
\cite[p.~45]{OlshanskiiSapir2003} and again in
\cite[Section~2.2(IX)]{HaettelOsajda2025}. The torsion-by-cyclic
groups of \cite[Theorem~1.1]{OlshanskiiSapir2003} have an infinite
cyclic quotient. They establish a different finite-presentation result
and are not themselves periodic. The companion paper
\cite{Companion} constructs an infinite finitely presented periodic
group using a graded algebra with a matrix-nil positive-degree ideal.
The present paper strengthens
that construction by proving residual finiteness and passing to a
finite-index subgroup with 2-power torsion. The companion's algebra
and finite local-rule hierarchy are essential inputs; we state their
precise properties in Section~\ref{sec:ring-phases}.

A bound on all element orders would change the problem. Novikov and
Adian constructed infinite finitely generated groups of bounded exponent
\cite{NovikovAdian1968}. In the residually finite setting, however,
Zel'manov's solution of the restricted Burnside problem for 2-groups
\cite[Theorem~2]{Zelmanov1992} implies that a finitely generated
residually finite group of exponent dividing $2^a$ is finite.
Consequently the group in Theorem~\ref{thm:main} has unbounded
exponent. Its finite quotients are nevertheless all 2-groups: a finite
image has only elements of 2-power order, and Cauchy's theorem excludes
every odd prime divisor of its order.

\subsection*{The proof and its main ingredients}

The companion supplies an infinite finitely presented graded algebra
\[
 R=\F_2\oplus I,\qquad I=\bigoplus_{d\ge1}R_d,
\]
such that every matrix in $M_n(I)$ is nilpotent, for every finite $n$.
Each homogeneous component $R_d$ is finite-dimensional. The elementary
matrix group $\E_n(R)$ is generated by the matrices $e_{ij}(a)$ obtained
by adding $a\in R$ to an off-diagonal position of the identity.
The Steinberg group $\St_n(R)$ is generated by symbols $x_{ij}(a)$
subject to the elementary addition and nonopposite-root commutator
relations; these are written explicitly in
Section~\ref{sec:steinberg}. There is a natural map
\[
 \St_n(R)\longrightarrow\E_n(R),\qquad
 x_{ij}(a)\longmapsto e_{ij}(a).
\]
Finite presentability of Steinberg groups in rank at least four over
finitely presented unital rings is due to Krsti\'c and McCool
\cite[Theorem~3]{KrsticMcCool}; an adapted proof is also given in
\cite[Theorem~4.1]{Companion}. We take $n=12$ and eventually use
the kernel of reduction from $\St_{12}(R)$ to $\St_{12}(\F_2)$.

The degree truncations $R/R_{\ge d}$, where
$R_{\ge d}=\bigoplus_{b\ge d}R_b$, are finite rings. They immediately
detect nonidentity matrices over $R$. The main difficulty is to detect
elements of the Steinberg kernel as well. We first work in the stable
group $\St(R)=\varinjlim_n\St_n(R)$, put
$\E(R)=\varinjlim_n\E_n(R)$, and write
\[
 J(R)=\ker\bigl(\St(R)\longrightarrow\E(R)\bigr).
\]
This kernel is central. We first prove detection in the stable group,
where a finite calculation may always use additional coordinates.
We return to fixed rank afterward.

Two calculations of one cyclic substitution provide the detection
argument. If $z\in J(R)$, the grading replaces each homogeneous root
entry $a\in R_d$ by $at^d$. Substituting the cyclic shift matrix $C_L$
for $t$ and expanding into ordinary roots gives an element
$W_z(C_L)\in J(R)$. This is a matrix realization of a graded transfer.
Weibel's work on graded rings \cite{Weibel1987} develops the
corresponding transfer operations and their eventual vanishing; it also
implies $p$-primary torsion in relative graded $K$-groups in
characteristic $p$. Following the elementary Steinberg treatment in
\cite{Companion}, Section~\ref{sec:shifts} proves directly that
\[
 W_z(C_L)=1\quad\text{for all sufficiently large }L
\]
whenever $z$ reduces to $1$ over $\F_2$. Taking $L$ to be a power of
two also shows that every such $z$ has 2-power order.

The second calculation uses additional structure of this particular
ring. Its finite local rules give a basis of nonzero word classes.
For arbitrarily large odd integers $L$, each sufficiently long nonzero
word has a start phase $\sigma(w)\in\Z/L$ and an end phase
$\tau(w)=\sigma(w)+|w|$. These phases depend only on the word class,
and a nonzero product must match end phase to start phase.
Section~\ref{sec:ring-phases} derives the phases from the hierarchy's
bounded token decompositions.

If $z$ is missed by all degree truncations, it can be represented using
roots whose entries are arbitrarily long words. The phases allow us
to move the endpoints of each such root to coordinates labeled by its
start and end phases, without changing $z$. Establishing this identity
requires control of opposite-root commutators in the Steinberg group,
as well as a subgroup on which the matrix map is injective.
Sections~\ref{sec:long-entries} and \ref{sec:detection} prove these
facts and obtain, for the long-entry representative,
\[
 W_z(C_L)=z^L.
\]
We first choose an odd $L$ for which the first calculation vanishes,
and then find a representative long enough for the second calculation.
Thus $z^L=1$. Since $z$ also has 2-power order, it is the identity.
This proves separation by finite degree truncations, the additional
step beyond torsion supplied by graded transfers alone.

To pass back to twelve coordinates, we need the injectivity of
$\St_{12}(R)\to\St(R)$. Stability over the local ring $R$ is
covered by the classical theorem of Suslin and Tulenbaev
\cite{SuslinTulenbaev1976}. Section~\ref{sec:stability} gives a
direct integral homology proof using ordered frames and universal
central extensions of their stabilizers. This supplies the bridge
from stable detection to the finitely presented group.

All arguments establishing these phase, stabilization, and detection
claims are included below. We quote from the companion the algebra
construction and the finite-rank elementary facts needed to apply them,
with their hypotheses stated at the points of use. The final passage
to finite groups and to the subgroup of 2-power torsion is carried out
in Section~\ref{sec:construction}.
